\chapter{Całka Lebesgue'a}
\label{chap:calka-lebesguea}

W tym rozdziale $(X,\mathcal{F},\mu)$ jest przestrzenią z miarą.
Schematy z tablicy przedstawiają zastępowanie funkcji prostokątami:
przez podział dziedziny oraz przez podział wartości funkcji
(rysunek~\ref{fig:calka-podzialy}).

\begin{figure}[htbp]
  \centering
  \includegraphics[width=0.96\textwidth]{calka-podzialy.pdf}
  \caption{Schematy podziału dziedziny i podziału wartości funkcji.}
  \label{fig:calka-podzialy}
\end{figure}

\section{Funkcje charakterystyczne}
\label{sec:funkcje-charakterystyczne}

Dla $A\subseteq X$ funkcję charakterystyczną zbioru $A$ oznaczamy
przez $\mathbf{1}_A$:
\[
  \mathbf{1}_A(x)=
  \begin{cases}
    1, & x\in A,\\
    0, & x\notin A.
  \end{cases}
\]

\begin{fakt}
\label{fakt:mierzalnosc-charakterystycznej}
Funkcja $\mathbf{1}_A$ jest mierzalna wtedy i tylko wtedy,
gdy zbiór $A$ jest mierzalny, czyli $A\in\mathcal{F}$.
\end{fakt}

\begin{proof}
Dla $a\in\R$ mamy
\[
  \{x:\mathbf{1}_A(x)>a\}
  =\begin{cases}
    X, & a<0,\\
    A, & 0\leq a<1,\\
    \varnothing, & a\geq1.
  \end{cases}
\]
Jeżeli $A\in\mathcal{F}$, wszystkie te zbiory są mierzalne.
Odwrotnie, jeżeli $\mathbf{1}_A$ jest mierzalna, to biorąc $a=0$,
otrzymujemy $A=\{x:\mathbf{1}_A(x)>0\}\in\mathcal{F}$.
\end{proof}

\section{Funkcje proste i postać kanoniczna}
\label{sec:funkcje-proste}

\begin{definicja}[Funkcja prosta]
\label{def:funkcja-prosta}
Funkcję $f\colon X\to\R$ nazywamy prostą, jeżeli przyjmuje
tylko skończenie wiele wartości.
\end{definicja}

\begin{definicja}[Postać kanoniczna]
\label{def:postac-kanoniczna}
Niech $f(X)=\{c_1,\ldots,c_r\}$, gdzie wartości $c_i$ są parami różne.
Połóżmy $E_i=\{x\in X:f(x)=c_i\}$. Wtedy
\begin{equation}
  f(x)=\sum_{i=1}^{r}c_i\mathbf{1}_{E_i}(x).
  \label{eq:postac-kanoniczna-prostej}
\end{equation}
Zapis ten nazywamy postacią kanoniczną funkcji prostej.
Zbiory $E_i$ są niepuste, parami rozłączne i
$\bigcup_{i=1}^{r}E_i=X$.
\end{definicja}

\begin{przyklad}[Dwa zapisy tej samej funkcji]
\label{prz:prosta-dwa-zapisy}
Funkcję z tablicy można zapisać jako
\[
  f=2\mathbf{1}_{[0,1]}+\mathbf{1}_{(1,2]}
    =\mathbf{1}_{[0,2]}+\mathbf{1}_{[0,1]}.
\]
Oba zapisy dają wartość $2$ na $[0,1]$, wartość $1$ na $(1,2]$
i wartość $0$ poza $[0,2]$.
W drugim zapisie zbiory $[0,2]$ i $[0,1]$ nie są rozłączne.
Pełna postać kanoniczna dla dziedziny $\R$ to
\[
  f=0\mathbf{1}_{\R\setminus[0,2]}
    +\mathbf{1}_{(1,2]}+2\mathbf{1}_{[0,1]}.
\]
Wykres przedstawia rysunek~\ref{fig:funkcje-proste} po prawej.
\end{przyklad}

\begin{figure}[htbp]
  \centering
  \includegraphics[width=0.96\textwidth]{funkcje-proste.pdf}
  \caption{Funkcja charakterystyczna przedziału oraz funkcja prosta z przykładu~\ref{prz:prosta-dwa-zapisy}.}
  \label{fig:funkcje-proste}
\end{figure}

\begin{twierdzenie}
\label{tw:mierzalnosc-prostej}
Jeżeli funkcja prosta $f$ ma postać kanoniczną
\eqref{eq:postac-kanoniczna-prostej}, to jest mierzalna wtedy i tylko
wtedy, gdy wszystkie zbiory $E_1,\ldots,E_r$ należą do $\mathcal{F}$.
\end{twierdzenie}

\begin{proof}
Jeżeli $f$ jest mierzalna, to
$E_i=f^{-1}(\{c_i\})\in\mathcal{F}$.
Odwrotnie, jeżeli wszystkie $E_i$ są mierzalne, to
funkcje $\mathbf{1}_{E_i}$ są mierzalne na mocy
faktu~\ref{fakt:mierzalnosc-charakterystycznej}.
Ich skończona kombinacja liniowa jest mierzalna.
\end{proof}

\section{Całka funkcji prostej: etap pierwszy}
\label{sec:calka-prostej-etap-pierwszy}

\begin{definicja}[Całka funkcji prostej]
\label{def:calka-prostej}
Dla mierzalnej funkcji prostej w postaci kanonicznej
\eqref{eq:postac-kanoniczna-prostej} definiujemy
\begin{equation}
  \int_X f\,d\mu=\sum_{i=1}^{r}c_i\mu(E_i),
  \label{eq:calka-prostej}
\end{equation}
o ile suma po prawej stronie jest określona.
W szczególności dla funkcji nieujemnej jest ona zawsze określona
i może być równa $+\infty$. Stosujemy konwencję $0\cdot(+\infty)=0$.
\end{definicja}

\begin{uwaga}
Jeżeli funkcja przyjmuje wartości obu znaków, zapis
\eqref{eq:calka-prostej} wymaga wykluczenia nieokreślonej sumy
$+\infty+(-\infty)$.
\end{uwaga}

\begin{wniosek}
\label{wn:calka-charakterystycznej}
Dla każdego zbioru $E\in\mathcal{F}$ zachodzi
\[
  \int_X\mathbf{1}_E\,d\mu=\mu(E).
\]
\end{wniosek}

\begin{proof}
W postaci kanonicznej występuje wartość $1$ na $E$ i wartość $0$
na $E^c$; puste poziomice pomijamy. Stąd całka jest równa
$1\cdot\mu(E)+0\cdot\mu(E^c)=\mu(E)$.
\end{proof}
